\documentclass[12pt]{article}

% --- Options passed, but may be ignored if paper loads biblatex with explicit options ---
\PassOptionsToPackage{backend=biber, style=numeric, sorting=none}{biblatex}

\usepackage{paper,math}

\usepackage[T2A]{fontenc}
\usepackage[utf8]{inputenc}
\usepackage[russian,english]{babel}
\usepackage{csquotes}

\addbibresource{references.bib}

\hfuzz=30pt % Suppresses warnings for lines up to 20pt too wide
\sloppy

\title{Title}
\author{
  Мельничук Михаил\\ {\itshape\large Университет ИТМО}
}
\date{\today}

\hypersetup{
  pdftitle = {Example Paper Title},
  pdfauthor = {Author One, Author Two},
}

\booltrue{INCLUDECOMMENTS}
\newcommand{\kyle}[1]{\coauthorComment[Kyle]{#1}}

\begin{document}

% Title Page -------------------------------------------------------------------
\maketitle
\begin{abstract}
  Abstract text

  \par\vspace{-2.5mm}
\end{abstract}

\newpage

% ##################################################################################################
\section{Introduction}

Citation example \cite{mas1995microeconomic}.

% ##################################################################################################
\section{Model}

Text text text text text text text text text text.

\graphicspath{{Melnichuk_Figs/}}
\begin{figure}[!htbp]
  \centering\small
  \begin{tabular}{cp{10mm}c}
    \includegraphics[height=1.3cm]{System_States.png} \\
  \end{tabular}\\
  \renewcommand{\figurename}{Fig.}
  \parbox{150mm}{
    \caption{\centering
      A transition diagram of the state space of the system. Each state $S_k$ corresponds to having $k$ requests in the system (either in service or waiting).
    }
    \label{Mel_fig_System_States}}
  \vspace{-1\baselineskip}
\end{figure}

Let $p_k$ denote steady state probability of state $S_k$. Solving global balance equations yields probability of having no requests  in the system $p_0$ is shown in Eq. \ref{eq:p0}.

\begin{align}
  p_0 & = \left(1+\sum_{i=1}^{n}\left(\frac{\lambda^i}{i!\prod_{j=1}^{i}\mu\left(n,\vec{R},j\right)}\right)+\sum_{r=1}^{m}\left(\frac{\lambda^{n+r}}{n!\prod_{j=1}^{n}\mu\left(n,\vec{R},j\right) n^r{\mu\left(n,\vec{R},n\right)}^r}\right)\right)^{-1} \nonumber                                                                                                            \\
      & = \left(1+\sum_{i=1}^{n}\left(\frac{\lambda^i}{i!\prod_{j=1}^{i}\mu\left(n,\vec{R},j\right)}\right)+\frac{\lambda^n}{n!\prod_{j=1}^{n}\mu\left(n,\vec{R},j\right)}\sum_{r=1}^{m}\left(\frac{\lambda}{n\mu\left(n,\vec{R},n\right)}\right)^r\right)^{-1} \nonumber                                                                                                     \\
      & = \left(1+\sum_{i=1}^{n}\left(\frac{\lambda^i}{i!\prod_{j=1}^{i}\mu\left(n,\vec{R},j\right)}\right)+\frac{\lambda^n}{n!\prod_{i=1}^{n}\mu\left(n,\vec{R},i\right)}\frac{\frac{\lambda}{n\mu\left(n,\vec{R},n\right)}\left(1-\left(\frac{\lambda}{n\mu\left(n,\vec{R},n\right)}\right)^m\right)}{1-\frac{\lambda}{n\mu\left(n,\vec{R},n\right)}}\right)^{-1} \nonumber \\
      & = \left(1+\sum_{i=1}^{n}\left(\frac{\lambda^i}{i!\prod_{j=1}^{i}\mu\left(n,\vec{R},j\right)}\right)+\frac{\lambda^{n+1}}{n!\prod_{i=1}^{n}\mu\left(n,\vec{R},i\right)}\frac{1-\left(\frac{\lambda}{n\mu\left(n,\vec{R},n\right)}\right)^m}{n\mu\left(n,\vec{R},n\right)-\lambda}\right)^{-1}
  \label{eq:p0}
\end{align}


\begin{equation}
  p_k =
  \begin{cases}
    \dfrac{\lambda^k}{k! \prod_{j=1}^{k} \mu(n,\vec{R},j)} \, p_0,                                         & k \le n       \\[10pt]
    \dfrac{\lambda^k}{\bigl(n \mu(n,\vec{R},n)\bigr)^{k-n} \, n! \prod_{i=1}^{n} \mu(n,\vec{R},i)} \, p_0, & n < k \le n+m \\[10pt]
    0,                                                                                                     & k > n+m
  \end{cases}
  \label{eq:pk}
\end{equation}

Среднее число заявок, находящихся под обслуживанием, т.е. среднее число занятых каналов:
\begin{equation}
  \begin{split}
    L_{\text{Об}} & = \sum_{r=1}^{n} r p_r + \sum_{r=1}^{m} n p_{n+r}                                                                                                                                                                                                    \\
                  & = \sum_{r=1}^{n} r \frac{\lambda^r}{r!\prod_{i=1}^{r}\mu(n,\vec{R},i)} p_0 + \sum_{r=1}^{m} n \frac{\lambda^{n+r}}{\bigl(n\mu(n,\vec{R},n)\bigr)^r \, n! \prod_{i=1}^{n}\mu(n,\vec{R},i)} p_0                                                        \\
                  & = \left( \sum_{r=1}^{n} \frac{\lambda^r}{(r-1)!\prod_{i=1}^{r}\mu(n,\vec{R},i)} + \frac{\lambda^n}{(n-1)!\prod_{i=1}^{n}\mu(n,\vec{R},i)} \sum_{r=1}^{m} \left(\frac{\lambda}{n\mu(n,\vec{R},n)}\right)^r \right) p_0                                \\
                  & = \left( \sum_{k=1}^{n} \frac{\lambda^k}{(k-1)!\prod_{i=1}^{k}\mu(n,\vec{R},i)} + \frac{\lambda^{n+1}}{(n-1)!\prod_{i=1}^{n}\mu(n,\vec{R},i)} \cdot \frac{1-\left(\frac{\lambda}{n\mu(n,\vec{R},n)}\right)^m}{n\mu(n,\vec{R},n)-\lambda} \right) p_0
  \end{split}
  \label{eq:L}
\end{equation}

% ##################################################################################################
\section{Results}

Text text text text text text text text text text.

% ##################################################################################################
\section{Results}

Text text text text text text text text text text.

% ##################################################################################################
\section{Discussion}
Text text text text text text text text text text.

% ##################################################################################################
\printbibliography[title={References}]   % <-- ADDED – prints the reference list

\end{document}